\section{空间周期与时间周期各自保存什么标记}
\label{chap:feqo-bloch-floquet}

若结构沿 \(z\) 每隔 \(a\) 重复，电子在 \(z\) 与 \(z+a\) 处看见相同势。定义平移算符 \((T_a\psi)(z)=\psi(z+a)\)。由 \(V(z+a)=V(z)\) 可直接核对 \([H,T_a]=0\)。先把体系放进长度为整数个 \(a\) 的大周期盒：此时能量子空间可按酉算符 \(T_a\) 的本征值分解，再取盒长趋于无穷的极限，所得无限结构的态可理解为广义本征态。因平移保持范数，\(T_a\) 的本征值模为一，写成 \(e^{\ii ka}\)。这就给出
\begin{equation}
\psi_k(z)=e^{\ii kz}u_k(z),\qquad u_k(z+a)=u_k(z).
\label{eq:feqo-bloch-state}
\end{equation}
证明只需令 \(u_k=e^{-\ii kz}\psi_k\)，再代入 \(\psi_k(z+a)=e^{\ii ka}\psi_k(z)\)。波数 \(k\) 与 \(k+G_m\)（\(G_m=2\pi m/a\)）对应相同平移本征值；周期函数相应乘 \(e^{-\ii G_mz}\)。这叫准动量，并不意味着电子真实动量只有一个值。

周期势的 Fourier 级数 \(V(z)=\sum_mV_me^{\ii G_mz}\) 立即显示矩阵元只连接 \(p\) 与 \(p+\hbar G_m\)。因此 Bragg 共振是这两个通道的自由能量近简并，再由 \(V_m\) 打开能隙；若 \(V_m=0\)，只有几何交点，不产生散射。

若 哈密顿量 随时间周期变化 \(H(t+T)=H(t)\)，普通定态 \(e^{-\ii Et/\hbar}|u\rangle\) 一般不再成立。定义一周期传播算符 \(U(T,0)\)。在有限维截断或它确有正规本征态的子空间内，取 \(U(T,0)|u_\alpha(0)\rangle=e^{-\ii\varepsilon_\alpha T/\hbar}|u_\alpha(0)\rangle\)；本征值模为一，因为传播算符酉。从该本征态出发，把精确解的相位 \(e^{-\ii\varepsilon_\alpha t/\hbar}\) 提出，余下部分每周期重复，得到
\begin{equation}
|\psi_\alpha(t)\rangle=e^{-\ii\varepsilon_\alpha t/\hbar}|u_\alpha(t)\rangle,
\qquad |u_\alpha(t+T)\rangle=|u_\alpha(t)\rangle.
\label{eq:feqo-floquet-state}
\end{equation}
\(\varepsilon\) 与 \(\varepsilon+m\hbar\Omega\) 给同一个一周期本征值，其中 \(\Omega=2\pi/T\)；这是准能的模 \(\hbar\Omega\) 自由度。式~\eqref{eq:feqo-floquet-state} 只要求周期性，不要求驱动很快。

上式的周期性可直接从传播律验证。周期驱动满足 \(U(t+T,T)=U(t,0)\)，所以 \(U(t+T,0)=U(t,0)U(T,0)\)。让它作用在选定的单周期本征态上，再乘 \(e^{\ii\varepsilon(t+T)/\hbar}\)，多出的 \(e^{-\ii\varepsilon T/\hbar}\) 与相位因子正好抵消，留下 \(|u(t+T)\rangle=|u(t)\rangle\)。若哈密顿量只有空间周期而没有时间周期，不能把这段证明借来谈准能；反过来时间周期也不会自动给电子真实动量一个倒格矢。

这里还有一个无限维的边界：酉算符可以有连续谱，未必有一组可数的正规本征向量张成整个空间。因此上式证明了每个存在的 Floquet 本征通道如何演化，并不宣称任意开放传播问题都能写成离散的 \(\sum_\alpha\)；连续谱部分应改用相应的谱积分。这个限定与 \(\varepsilon\) 模 \(\hbar\Omega\) 的记号自由度是两回事。

高频有效 哈密顿量 是额外近似。若非零 Fourier 分量 \(H_m\) 的相关矩阵元满足 \(\|H_m\|/(\hbar\Omega)\ll1\)，远离多光子共振，便可消去快速振荡到指定阶；被消去的周期内运动仍以微运动算符存在。若驱动频率接近两个能级差的整数分之一，即使 \(\Omega\) 数值很大，这个展开也可能因为小分母而失效。

\begin{exercise}
直接代入验证 \(k\to k+G_m\)、\(u_k\to e^{-\ii G_mz}u_k\) 不改变 Bloch 态；再验证 \(\varepsilon\to\varepsilon+m\hbar\Omega\) 时怎样重定义周期态。
\end{exercise}
